\begin{abstract}
We give fully distributed multi-point function (DMPF) constructions that
share a sparse map over a domain of size $N$.
Setup takes secret-shared indices; each later expansion takes new payloads
for those indices and returns shares of the map.
Our main construction, Waterfall Cuckoo, samples public hashes before
assigning the inputs to bins. An occupancy scan followed by bounded repair
finds the assignment; a hidden permutation conceals it during routing.
Sparse DPFs then provide $4N$ or $3N$ leaf work per expansion, replacing the
$tN$ work of $t$ independent DPFs.
We prove correctness and semi-honest simulation in the ideal-subprotocol
model, with an explicit routing-error bound.

For the random sparse supports in Ring-LPN PCGs, we also give Reverse Cuckoo,
a $2N$-expansion specialization. It programs hashes around a hidden placement
and leaks information about the support: in the ideal hash model, outputs
are weighted by their number of valid placements.
We characterize the idealized leakage, state the additional decisional
Ring-LPN assumption, and evaluate leakage-aware decoding attacks.
A separate analysis bounds the concrete hash evaluator's batch correctness
error below $2^{-40}$ at the filtered degree-$2^{20}$ parameter point.

Our measurements use an earlier implementation of this specialized
Reverse-Cuckoo pipeline. They predate the mask, support-filter, permutation,
and evaluator-seed updates, and the AES rekeying described here.
At degree $2^{20}$ it produces 1.81 million Goldilocks-field OLEs per second,
with 12.99 MB communicated per expansion.
Against the reported FHE baseline at the same field and degree, these figures
give $6.57\times$ higher throughput and $13.60\times$ lower expansion
communication. Setup costs and comparisons with the sum-of-DPFs baseline
are reported separately.
Waterfall's costs are modeled; gains for other fields, rings, and PCG
constructions remain application-dependent.
\end{abstract}
